% GENERATED FILE. Edit canonical lessons, then run npm run generate:books.
% canonical-source-hash: fnv1a64:39c3d8b2984502df
% canonical-lessons: ES-C61-al

\chapter{Al --- Two Words Spanish Will Not Keep Apart}
\label{ch:al}
\hlchaptermodality{230}

\begin{chapteropening}
\textbf{By the end of this chapter:} \emph{I can say where I am going, with the article fused on.}

I can say where I am going, with the article fused on.
\end{chapteropening}

\section[al]{al — two words Spanish will not let you keep apart}

\label{lesson:ES-C61-al}



\begin{quote}
You have \emph{a} — \emph{to}, \emph{towards}. And you have \emph{el}.

Put them together and say it quickly. \emph{A el}... \emph{a el}... Hear what your mouth wants to do.
\end{quote}

\begin{grammarlens}[title={Grammar Lens: the squeeze}]
\begin{quote}
\emph{a} + \emph{el} $\to$ \emph{\textbf{al}}
\end{quote}

\begin{quote}
\emph{Voy} \emph{\textbf{al}} \emph{café.} — I'm going to the café.
\end{quote}

This is not an option and not a shortcut for speaking fast. It is how the language is \textbf{written}. \emph{A el café} is not a more careful version of \emph{al café}; it is simply wrong.
\end{grammarlens}

\begin{culture}
You met this instinct two chapters ago, when \emph{mucho} lost its ending in front of an adjective. Spanish does not like a word standing about doing nothing while another word waits: \textbf{when one word leans on the next, something gives.} Here the two words simply fuse.

And there is one exception worth knowing, because it looks like a mistake and is not. When \emph{El} is part of a \textbf{name}, it keeps its distance:

\begin{quote}
\emph{Voy a El Salvador.} — not \emph{al Salvador}.
\end{quote}

What contracts is the little article — not every word that happens to be spelled \emph{el}.
\end{culture}

\subsection*{Your turn}
Say these aloud:

\begin{itemize}
\raggedright
  \item \emph{al café}, \emph{al profesor}
  \item \emph{Voy al café}
\end{itemize}

\emph{Twice through:}

\begin{tcolorbox}[breakable,colback=teal!4,colframe=teal!35!black,title={Before you move on}]
\emph{A} plus \emph{el} makes what? (\emph{\textbf{Al}} — always.)

And when does it not happen? (When \emph{\textbf{El}} is part of a \textbf{name}.)
\end{tcolorbox}
